\section{Constrained Parameter Geometry and Asymptotics}
\label{sec:geometry_asymptotics}

\subsection{Geometry of the covariance parameter space}
\label{subsec:covariance_geometry}
Let $\lambda\in\Lambda\subset\mathbb R^{p_\lambda}$ collect the nuisance parameters, and let $\gamma=(\gamma_{11},\gamma_{12},\gamma_{22})^{\mathsf T}$. For the random-intercept-and-slope model, write
\begin{equation}\label{eq:sec4-psd-cone}
\Gamma(\gamma)=
\begin{pmatrix}
\gamma_{11}&\gamma_{12}\\
\gamma_{12}&\gamma_{22}
\end{pmatrix}\succeq0.
\end{equation}
The covariance constraint and full parameter space are
\begin{equation}\label{eq:sec4-parameter-space}
\begin{aligned}
K&=\{\gamma\in\mathbb R^3:\Gamma(\gamma)\succeq0\},\\
\Theta&=\Lambda\times K.
\end{aligned}
\end{equation}
We assume throughout that the true nuisance parameter $\lambda_0$ is interior to $\Lambda$. The covariance matrix $\Gamma_0$ may be positive definite, rank one, or zero. These cases have different feasible local perturbations and therefore different asymptotic behavior.

\subsection{Lorentz-cone representation of the $2\times2$ PSD cone}
\label{subsec:lorentz_representation}
Define $t=L\gamma$ by
\begin{equation}\label{eq:sec4-lorentz-map}
\begin{aligned}
t_0&=\tfrac12(\gamma_{11}+\gamma_{22}),\\
t_1&=\tfrac12(\gamma_{11}-\gamma_{22}),\\
t_2&=\gamma_{12}.
\end{aligned}
\end{equation}
This is an invertible linear transformation, under which $\Gamma\succeq0$ is equivalent to
\begin{equation}\label{eq:sec4-lorentz-cone}
\begin{aligned}
t\in\mathcal L_3
=\bigl\{t\in\mathbb R^3:\;&t_0\geq0,\\
&t_0^2\geq t_1^2+t_2^2\bigr\}.
\end{aligned}
\end{equation}
Thus $LK=\mathcal L_3$, the three-dimensional Lorentz cone. Its vertex, boundary rays, and interior correspond to zero, rank-one, and positive-definite covariance matrices. This is a coordinate representation of the PSD constraint; it does not make the Fisher information Euclidean. The matrix $L$ and the corresponding geometric transformations are given in Appendix~C.


\subsection{Tangent cones at interior, rank-one,
and vertex points}
\label{subsec:tangent_cones}

For a closed set $C$ and $c\in C$, let
$T_C(c)$ denote the Bouligand (contingent)
tangent cone. Tangent cones describe the feasible first-order
perturbations of a constrained parameter space.
Their use in boundary likelihood asymptotics is
standard; see \citet{self1987asymptotic} and
\citet{shapiro1985}.
The tangent cones of the PSD cone in this
setting are closed and convex.


\begin{proposition}\label{prop:sec4-tangent-cones}
Let $\Gamma_0\succeq0$, with covariance coordinates $\gamma_0$. Then:
\begin{enumerate}
\item If $\Gamma_0\succ0$, then $T_K(\gamma_0)=\mathbb R^3$.
\item If $\operatorname{rank}(\Gamma_0)=1$, let $u_0$ be a unit vector spanning $\ker(\Gamma_0)$. Then
\begin{equation}\label{eq:sec4-rankone-tangent}
\begin{aligned}
T_K(\gamma_0)=\{h_\gamma:\;&u_0^{\mathsf T}
H(h_\gamma)u_0\geq0\},
\end{aligned}
\end{equation}
where
\[
H(h_\gamma)=
\begin{pmatrix}
h_{11}&h_{12}\\
h_{12}&h_{22}
\end{pmatrix}.
\]
\item If $\Gamma_0=0$, then $T_K(0)=K$.
\end{enumerate}
Moreover, because $\lambda_0$ is interior,
\begin{equation}\label{eq:sec4-full-tangent}
T_\Theta(\theta_0)
=\mathbb R^{p_\lambda}\times T_K(\gamma_0).
\end{equation}
\end{proposition}
At a rank-one covariance matrix, the tangent cone is a half-space, although the feasible covariance matrices form a curved cone. A tangent direction need not itself be PSD: it need only arise as the first-order limit of feasible perturbations around $\Gamma_0$. Appendix~C proves the characterization, including its Lorentz-coordinate form.

\subsection{Local asymptotic normality and quadratic likelihood expansion}
\label{subsec:lan}
Let $\ell_n(\theta)$ denote the log-likelihood for the full parameter $\theta=(\lambda^{\mathsf T},\gamma^{\mathsf T})^{\mathsf T}$, and let $\theta_0\in\Theta$ be the true parameter. Let $I_0$ be the positive-definite limiting information matrix for the full parameter vector. Following the general boundary-likelihood framework of \citet{self1987asymptotic}, assume the following local asymptotic quadratic expansion holds uniformly over bounded feasible local perturbations $h$:
\begin{equation}\label{eq:sec4-lan}
\begin{aligned}
&\ell_n(\theta_0+h/\sqrt n)-\ell_n(\theta_0)\\
&\quad=h^{\mathsf T}Z_n
-\tfrac12h^{\mathsf T}I_0h+o_p(1),
\end{aligned}
\end{equation}
where $Z_n\xrightarrow{d}Z$ and $Z\sim N(0,I_0)$. The expansion is restricted to $\sqrt n(\Theta-\theta_0)$; it does not require likelihood evaluation at an indefinite covariance matrix.


We further assume local identifiability, consistency,
root-$n$ tightness, and sufficiently accurate likelihood
maximization for the constrained estimators under
consideration. The local feasible parameter sets are
assumed to converge to their corresponding tangent
cones in the sense required by the constrained argmax
theorem.

For likelihood-ratio inference, these conditions
apply separately to the full-model and restricted-null
maximum-likelihood estimators. In particular, both
estimators must be localized at the root-$n$ scale,
and their likelihood optimization errors must be
$o_p(1)$ on the local log-likelihood scale.
Together with the common uniform local quadratic
expansion, these conditions ensure joint convergence
of the maximized likelihood criteria.

These are high-level assumptions on the likelihood
and its maximizers, not claims of automatic
regularity for any particular response model.
The results below apply to likelihoods satisfying
these conditions, independently of their specific
conditional response distributions.
Appendix~\ref{app:asymptotics} develops the
corresponding constrained quadratic optimization
and restricted-null limits.

\subsection{Scope of compact-uniform inference}
\label{subsec:sec4-compact-uniform}

The preceding likelihood conditions are stated at a fixed true
parameter. For the compact-uniform results developed in subsequent
sections, we use a specified class of independent-cluster experiments
and a common enlarged compact fitting domain. Let
$\Lambda^+\subset\mathbb R^{p_\lambda}$ be a compact convex set
contained in the admissible nuisance-parameter domain, and let
$\Lambda_0\subset\operatorname{int}(\Lambda^+)$ be compact. For
$0<M_0<M<\infty$, define
\begin{equation}\label{eq:sec4-compact-classes}
\begin{aligned}
\Theta^+ & =\Lambda^+\times
 \{\gamma\in K:\operatorname{tr}\Gamma(\gamma)\leq M\},\\
\mathcal C & =\Lambda_0\times
 \{\gamma\in K:\operatorname{tr}\Gamma(\gamma)\leq M_0\}.
\end{aligned}
\end{equation}
The true parameter ranges over $\mathcal C$, whereas the unrestricted
and null-constrained likelihoods are maximized over the corresponding
subsets of $\Theta^+$. The gaps between these sets make the artificial
nuisance and upper-covariance boundaries locally inactive. Crucially,
$\mathcal C$ contains the covariance vertex, rank-one matrices, and
positive-definite matrices; no positive lower bound is imposed on the
eigenvalues of $\Gamma$.

The verified four-family applications use a fixed finite collection of
bounded cluster-design types with convergent frequencies, including
positive-frequency identifying designs. Fixed-effect and family-specific
nuisance dimensions are arbitrary but fixed. Gaussian residual variance
and NB2 size are bounded away from zero and infinity. In the compound
Poisson--Gamma family, both dispersion and power are unknown, with
dispersion in a compact positive interval and power in a compact
subinterval of $(1,2)$. The precise replicated-design and identifying
conditions are stated in the family-specific proofs.

Uniform conclusions over $\mathcal C$ additionally require
true-parameter-centered quadratic likelihood expansions, uniform score
normality and positive-definite information, and uniform localization of
\emph{both} exact global likelihood maximizers over $\Theta^+$ and its
relevant null subsets. These properties are verified on the specified
informative compact classes using the family-specific analytic results
and the general likelihood arguments of Appendix~\ref{app:general-regularity};
they are not asserted for arbitrary designs or unbounded fitting
domains. Replacing exact likelihood suprema by numerical fits requires
separate control of the entire maximized log-likelihood error on the
$o_p(1)$ scale. Thus fixed-parameter results below retain their original
high-level assumptions, whereas claims of compact-uniform size or
coverage use the stronger, explicitly delimited class above.

\subsection{Constrained-MLE asymptotic distribution}
\label{subsec:constrained_mle_limit}

The Gaussian projection characterization of
constrained likelihood estimators is a standard
consequence of boundary likelihood theory
\citep{self1987asymptotic,shapiro1985}.
We state its specialization to the present
PSD-constrained parameter space.
For a closed convex cone $T$ and a
positive-definite matrix $I$, define the
Fisher-metric projection
\begin{equation}
\Pi_T^I(y)
=
\underset{h\in T}{\operatorname{argmin}}
\,(h-y)^{\mathsf T}I(h-y).
\label{eq:sec4-metric-projection}
\end{equation}
Write $T_0=T_\Theta(\theta_0)$.


\begin{theorem}\label{thm:sec4-full-mle}
Suppose the conditions in Subsection~\ref{subsec:lan} hold at $\theta_0=(\lambda_0^{\mathsf T},\gamma_0^{\mathsf T})^{\mathsf T}$, with $\lambda_0$ interior and $\Gamma_0\succeq0$. Let $\widehat\theta_n$ be a consistent constrained maximum likelihood estimator. Then
\begin{equation}\label{eq:sec4-full-mle-limit}
\begin{aligned}
\sqrt n(\widehat\theta_n-\theta_0)
&\xrightarrow{d}H_0,\\
H_0&=\Pi_{T_0}^{I_0}(I_0^{-1}Z),
\end{aligned}
\end{equation}
where $Z\sim N(0,I_0)$. The projection acts on the full parameter vector, not only its covariance coordinates.
\end{theorem}
If $\Gamma_0\succ0$, then $T_0$ is Euclidean space and $H_0\sim N(0,I_0^{-1})$. At rank-one and vertex covariance parameters, the projection is generally nonlinear, and the full-vector limit is generally non-Gaussian. The tangent cone in the theorem is evaluated at the \emph{true} parameter $\theta_0$, not at $\widehat\theta_n$.


\begin{corollary}
\label{cor:sec4-restricted-null}
Suppose $\Gamma_0=0$ and the conditions of
Subsection~\ref{subsec:lan} hold for both
the full and restricted likelihood estimators.

Let $\widetilde\lambda_n$ denote the
maximum-likelihood estimator under
$H_0:\Gamma=0$. Then
\begin{equation}
\label{eq:sec4-restricted-null-limit}
\sqrt n(\widetilde\lambda_n-\lambda_0)
\xrightarrow{d}
N_{p_\lambda}
\left(0,I_{\lambda\lambda,0}^{-1}\right).
\end{equation}
Moreover, the maximized local likelihood
criteria under the full and restricted
parameter spaces converge jointly to
their respective Gaussian quadratic
suprema.
\end{corollary}

The result follows by maximizing the common local
quadratic criterion over the restricted tangent space
$\mathbb R^{p_\lambda}\times\{0\}$.
The maximizer is $I_{\lambda\lambda,0}^{-1}Z_\lambda$;
joint convergence follows from the common LAN process
and the localization assumptions for both fits.
The argument is detailed in Appendix~\ref{app:asymptotics}.


\subsection{Projection under the Fisher information metric}
\label{subsec:fisher_projection}
Partition the information matrix as
\begin{equation}\label{eq:sec4-information-partition}
I_0=
\begin{pmatrix}
I_{\lambda\lambda}&I_{\lambda\gamma}\\
I_{\gamma\lambda}&I_{\gamma\gamma}
\end{pmatrix},
\end{equation}
with blocks evaluated at $\theta_0$. Define the efficient covariance information
\begin{equation}\label{eq:sec4-efficient-information}
\begin{aligned}
J={}&I_{\gamma\gamma}\\
&-I_{\gamma\lambda}I_{\lambda\lambda}^{-1}
I_{\lambda\gamma}.
\end{aligned}
\end{equation}
Partition $Z=(Z_\lambda^{\mathsf T},Z_\gamma^{\mathsf T})^{\mathsf T}$ and set
\begin{equation}\label{eq:sec4-efficient-score}
\begin{aligned}
Z_{\gamma\cdot\lambda}
={}&Z_\gamma\\
&-I_{\gamma\lambda}I_{\lambda\lambda}^{-1}Z_\lambda.
\end{aligned}
\end{equation}
Then $Z_{\gamma\cdot\lambda}\sim N(0,J)$. Writing $H_0=(h_\lambda^{*\mathsf T},h_\gamma^{*\mathsf T})^{\mathsf T}$, the full projection admits the block characterization
\begin{align}
 h_\gamma^*
 &=\Pi_{T_K(\gamma_0)}^J
 (J^{-1}Z_{\gamma\cdot\lambda}),
 \label{eq:sec4-cov-projection}\\
 h_\lambda^*
 &=I_{\lambda\lambda}^{-1}
 (Z_\lambda-I_{\lambda\gamma}h_\gamma^*).
 \label{eq:sec4-nuisance-projection}
\end{align}
Thus nuisance adjustment determines the covariance projection through $J$, while the covariance constraint feeds back into the nuisance estimator through $I_{\lambda\gamma}$. Unrestricted nuisance coordinates do not, by themselves, imply an unconstrained Gaussian limit for their estimators. These block identities follow by profiling the Gaussian quadratic criterion; the calculation is given in Appendix~\ref{app:asymptotics}.


The efficient-information formula above is expressed in the original
covariance coordinates. When an exact nonlinear VPC hypothesis is
straightened by a nuisance-dependent transformation of covariance
coordinates, the information must first be transformed for the
\emph{full} parameter vector. If $F$ maps transformed coordinates
$\eta=(\lambda^{\mathsf T},\nu^{\mathsf T})^{\mathsf T}$ to the
original $\theta$, with Jacobian $L_F=D_\eta F$, then the appropriate
information is $\widetilde I=L_F^{\mathsf T}I L_F$ and its efficient
covariance block is
\begin{equation}\label{eq:sec4-transformed-efficient-information}
\widetilde J=\widetilde I_{\nu\nu}
-\widetilde I_{\nu\lambda}
 \widetilde I_{\lambda\lambda}^{-1}
 \widetilde I_{\lambda\nu}.
\end{equation}
A covariance-only congruence of $J$ is generally incorrect when the
coordinate chart depends on nuisance parameters; the distinction is
used for the exact VPC contrast geometry in
Section~\ref{sec:equality_vpc}.

\subsection{Implications for inference on smooth functionals}
\label{subsec:functional_asymptotics}

Let $g:\Theta\to\mathbb R^r$ denote a functional,
such as the VPC at one or more covariate values
or a contrast of VPCs. Suppose $g$ admits a
continuously differentiable local extension
around $\theta_0$, or more generally is
Hadamard directionally differentiable
tangentially to $T_0$, with derivative
$g'_{\theta_0}(h)$.
Applying the directional delta method of
\citet{fangsantos2019}, under its requisite
tangential differentiability and regularity
conditions, gives
\begin{equation}
\begin{aligned}
&\sqrt n\{g(\widehat\theta_n)-g(\theta_0)\}\\
&\qquad\xrightarrow{d}g'_{\theta_0}(H_0).
\end{aligned}
\label{eq:sec4-functional-limit}
\end{equation}
For a smooth extension, $g'_{\theta_0}(h)=Dg(\theta_0)h$. Even then, the limit need not be Gaussian because $H_0$ need not be Gaussian. At a zero-VPC point $x_0$, Proposition~\ref{prop:zero_vpc_gradient} identifies the ambient gradient, but feasible directions and the true covariance rank remain essential. If the first-order derivative vanishes identically on the limiting directions, higher-order analysis is required.


Section~\ref{sec:global_lrt} specializes the
vertex geometry to global likelihood-ratio
inference. Other inferential problems may
require different null tangent sets, while
the general projection framework remains the same.
